\documentclass{article}
% Use preprint mode for a readable, non-anonymous working draft.
% Switch to \usepackage{neurips_2025} for anonymous submission and
% \usepackage[main, final]{neurips_2025} only after acceptance.
\usepackage[preprint]{neurips_2025}

\usepackage[utf8]{inputenc}
\usepackage[T1]{fontenc}
\usepackage{hyperref}
\usepackage{url}
\usepackage{microtype}
\usepackage{xcolor}
\usepackage{amsmath}
\usepackage{booktabs}
\usepackage{enumitem}

\setlist[itemize]{leftmargin=*,itemsep=2pt,topsep=2pt,parsep=0pt,partopsep=0pt}

\title{Adaptive Minimal Evidence Acquisition for Grounded VideoQA \\ \large Literature, Positioning, and Implementation Status}

\author{
  Ritesh Thawkar \\
  Mohamed Bin Zayed University of Artificial Intelligence \\
  \texttt{ritesh.thawkar@mbzuai.ac.ae}
}

\begin{document}
\maketitle

\begin{abstract}
This document is a living research draft for the CV7502 project and a starting point for a potential conference submission.
It records the current literature review, the observed research landscape, the refined novelty claim, and the implementation status of the working codebase.
The central lesson from the literature is that multimodal VideoQA has already moved beyond simple fusion and generic retrieval; the active frontier is now grounded evidence selection, adaptive temporal retrieval, compute-aware context construction, and faithfulness of the retrieved evidence.
In response, the project has shifted from a broad tri-modal routing idea toward adaptive minimal evidence acquisition: under a fixed evidence budget, the system should retrieve only the smallest subtitle-frame-segment evidence set needed to answer correctly and faithfully.
\end{abstract}

\section{Project Framing}
The current project studies Video Question Answering (VideoQA) with evidence drawn from subtitles, frames, and short temporal segments under a fixed retrieval budget.
The original proposal argued that question-conditioned routing of the evidence budget across modalities could improve answer accuracy and reduce evidence mismatch relative to single-modality or fixed-split retrieval.

The literature review below suggests that this framing is directionally correct, but that modality routing alone is no longer sufficient as a publishable claim.
Recent work already studies question-guided temporal selection, long-video retrieval, grounded answer prediction, and modality-aware retrieval.
To produce a strong research contribution, the project will likely need to focus on \emph{budget-aware}, \emph{faithful}, and \emph{robust} evidence allocation rather than only on multi-modal retrieval.

\section{Objective and Intended Contribution}
The working objective is now:
\begin{quote}
\emph{Formulate grounded VideoQA as adaptive multimodal evidence acquisition, where the system sequentially chooses whether to acquire subtitle, frame, or segment evidence, or to stop, under a fixed evidence budget.}
\end{quote}

The intended contribution is not simply to use multiple modalities.
That contribution is already well represented in the literature.
The project instead aims to contribute the following combination:
\begin{itemize}
  \item \textbf{Budget-aware acquisition:} learn how much evidence to acquire and when to stop under a hard budget.
  \item \textbf{Minimal sufficient evidence:} build and imitate evidence sets that are as small as possible while preserving correctness.
  \item \textbf{Faithfulness constraints:} require the selected evidence to remain grounded in the correct temporal region and, when possible, retain answer support.
  \item \textbf{Reproducible evaluation:} report accuracy, cost, sufficiency, comprehensiveness, and temporal grounding together rather than accuracy alone.
\end{itemize}

\section{Literature Review}
\subsection{Foundational VideoQA benchmarks and supervised multimodal reasoning}
Early work established the task formulations and benchmark culture for VideoQA.
\citet{tapaswi2016movieqa} introduced MovieQA as a story-understanding benchmark for movies, while \citet{jang2017tgifqa} proposed TGIF-QA to push beyond static-image reasoning toward spatio-temporal reasoning in GIF videos.
The field consolidated around multi-modal reasoning with \citet{lei2018tvqa}, who introduced TVQA with aligned subtitles and localized temporal reasoning.
TVQA was particularly influential because it made subtitle evidence, question-conditioned localization, and multiple-choice answering part of a single benchmark.

Subsequent work pushed from answer prediction toward grounded evidence prediction.
\citet{lei2019tvqaplus} extended TVQA into TVQA+, adding spatio-temporal grounding annotations and proposing STAGE for joint answering and grounding.
Architectures such as HCRN \citep{le2020hcrn} and other hierarchical video-language fusion models showed that stronger temporal compositional structure could improve reasoning, but these systems still typically consumed dense video context rather than explicitly managing an evidence budget.
Datasets such as NExT-QA \citep{xiao2021nextqa} further exposed the limits of shallow scene-description methods by emphasizing temporal and causal reasoning.

The main takeaway from this first phase of the literature is that VideoQA is not just a classification problem.
High performance requires identifying which time span, which visual content, and which textual context actually support the answer.
This observation directly motivates retrieval-style formulations.

\subsection{Large-scale pretraining and frozen-language-model approaches}
The second phase of the literature shifted from dataset-specific fusion architectures to large-scale pretraining and frozen language models.
Representation learning advances such as CLIP \citep{radford2021clip}, VideoCLIP \citep{xu2021videoclip}, and CLIP4Clip \citep{luo2021clip4clip} made dense cross-modal retrieval practical for text-image and text-video matching.
These models became the backbone for many later retrieval and grounding pipelines.

On the QA side, \citet{yang2021justask} showed that large-scale automatically generated supervision from narrated videos can support strong VideoQA performance.
This line of work reduced the dependence on manually annotated QA pairs and made zero-shot and few-shot VideoQA more realistic.
\citet{yang2022frozenbilm} then showed that frozen bidirectional language models can be adapted for zero-shot VideoQA using lightweight trainable modules, providing a stronger alternative to heavier autoregressive approaches.
\citet{pan2023r2a} took this further with a retrieval-based design that converts video understanding into language-space reasoning by retrieving semantically similar text from an external corpus before answer prediction.

These methods revealed an important pattern: strong VideoQA performance can come from combining pretrained visual encoders, language-space alignment, and lightweight answer modules, without expensive end-to-end video training.
However, they also expose a limitation relevant to the current project.
Much of the ``retrieval'' in this phase is effectively text-side support retrieval or language-space prompting, not fine-grained evidence allocation across subtitles, frames, and temporal segments under a shared budget.

\subsection{Retrieval augmentation, long-video memory, and evidence compression}
Recent work has expanded retrieval from a supportive component into the core mechanism for video understanding, especially for long videos where full-context processing is too expensive.
\citet{luo2024videorag} propose Video-RAG for visually aligned retrieval-augmented long-video comprehension, making explicit the need to retrieve only the most relevant evidence rather than process entire videos.
Related systems such as VRAG \citep{gia2025vrag} and VideoRAG over video corpora \citep{jeong2025videorag} further push retrieval-augmented reasoning toward long-form VideoQA and corpus-scale retrieval settings.

This theme overlaps with long-video memory work.
Although not always cast as retrieval, the central problem is similar: a system must compress, select, or revisit only the relevant pieces of a long and redundant video stream.
The modern question is therefore not whether some retrieval mechanism helps, but what information should be retrieved, at what granularity, and under what compute budget.

For the current project, this literature matters because it strengthens the case for a fixed evidence-budget setup.
If the project compares methods under equal budgets, it can make a sharper claim about \emph{efficiency of evidence selection}, not just raw accuracy.
That framing is stronger both scientifically and for potential publication.

\subsection{Grounded and trustworthy VideoQA}
Another major recent direction is faithfulness and grounding.
\citet{xiao2023trust} show that strong VideoQA models may answer correctly while failing to ground their predictions in the relevant video evidence.
Their NExT-GQA benchmark and analysis are important because they separate answer correctness from answer trustworthiness.
This is directly relevant to any retrieval-augmented system: retrieving evidence is only useful if the answer is actually based on that evidence.

Grounded VideoQA work suggests that future systems should not only output an answer but also expose the supporting temporal span, frames, or subtitle spans.
For the present project, this creates a natural opportunity.
If the model returns cited subtitle chunks, timestamps, and frames, then routing can be evaluated not just as a hidden controller but as a mechanism for \emph{faithful evidence construction}.
This is likely a stronger contribution than a pure accuracy improvement.

\subsection{Question-guided temporal retrieval and modality-aware routing}
The most relevant recent literature is the work that directly adapts retrieval to the question.
\citet{amoroso2024pqr} propose question-guided temporal queries for VideoQA, showing that temporal evidence should be selected conditionally on the input question rather than extracted with fixed pooling.
This narrows the novelty space for any project whose main claim is only ``question-guided evidence selection.''

The closest line to the current proposal is modality-aware retrieval.
\citet{yeo2025universalrag} argue that retrieval over multiple corpora with diverse modalities and granularities suffers from modality gaps in a unified retrieval space, motivating modality- and granularity-aware retrieval.
\citet{delarosa2025modaroute} push this further with smart routing for multimodal video retrieval, where the system decides which modalities to search.
Together, these papers imply that a publishable contribution can no longer stop at saying that different questions prefer different modalities.
That point is already established.

The remaining open problem is finer and more operational: how should a system allocate a limited evidence budget across modalities and granularities so that retrieval is accurate, efficient, and faithful?
This budget-allocation perspective is not yet cleanly resolved in the current literature, especially in benchmark settings where subtitle and visual evidence compete.

\section{Synthesis of the Current Landscape}
The literature supports five clear conclusions.

\begin{itemize}
  \item VideoQA has evolved from late-fusion answer classification to evidence-centric reasoning with temporal grounding and explicit context selection.
  \item Pretrained vision-language encoders and frozen language models have made retrieval-based and lightweight-answering pipelines highly competitive.
  \item Long-video understanding has turned retrieval, memory compression, and context budgeting into central system design problems.
  \item Current models often remain weak in faithfulness: correct answers do not necessarily imply correct evidence.
  \item Modality-aware and question-conditioned retrieval already exist, so novelty now depends on how evidence is allocated, justified, and evaluated.
\end{itemize}

This means that the project's baseline idea remains viable, but it needs sharper positioning.
``Tri-modal routing'' by itself is not sufficiently new in 2026.
A stronger contribution would be one of the following:

\begin{itemize}
  \item \textbf{Budget-aware routing:} learn both \emph{which} modality matters and \emph{how much} budget each modality should receive under a hard evidence cap.
  \item \textbf{Faithful routed VideoQA:} require cited evidence and evaluate whether the routed evidence genuinely supports the predicted answer.
  \item \textbf{Robust routing:} show that modality routing adapts under noisy subtitles, missing modalities, or lower evidence budgets.
  \item \textbf{Hierarchical routing:} route across modality first and temporal granularity second, instead of using a flat top-$K$ retrieval rule.
\end{itemize}

\section{Working Research Position}
Based on the current survey, the strongest working direction is:
\begin{quote}
\emph{Study grounded VideoQA as a model-relative minimal-evidence problem: under a fixed budget, compare how weak and stronger answerers differ in the smallest subtitle-frame-segment evidence sets they need, and learn acquisition policies that are evaluated on accuracy, faithfulness, compactness, and transfer.}
\end{quote}

This positioning is more defensible than a generic tri-modal RAG claim because it ties the contribution to a concrete unresolved question: whether ``minimal sufficient evidence'' is stable across answerers or is itself model-dependent.
It also aligns well with course expectations around novelty, thorough experimentation, and critical analysis.

\section{Current Implementation Status}
The codebase now supports the main research pipeline at a baseline level.
Implemented components include:
\begin{itemize}
  \item normalization for TVQA and TVQA+ into a shared JSONL schema;
  \item normalization for NExT-GQA from official CSV annotations, grounding JSON, frame-time JSON, and video-id maps;
  \item candidate-pool construction for subtitle chunks, frame timestamps, and temporal segments;
  \item visual evidence materialization from source videos using \texttt{ffmpeg};
  \item CLIP-based feature extraction for frames and segment-level pooled embeddings;
  \item retrieval baselines including lexical, BM25, and hybrid subtitle-plus-CLIP retrieval;
  \item a fixed-budget linear answerer baseline and a stronger frozen multimodal answerer backend using CLIP text embeddings plus the visual feature store;
  \item a constrained greedy oracle for minimal evidence subsets, with explicit modes for prediction-preserving, correctness-only, correctness-plus-sufficiency, and correctness-plus-sufficiency-plus-grounding comparisons;
  \item oracle-trace export for imitation learning;
  \item a trainable sequential policy that predicts \texttt{acquire\_subtitle}, \texttt{acquire\_frame}, \texttt{acquire\_segment}, or \texttt{stop};
  \item a model-relative study runner that compares minimal-evidence subsets across two answerers and records evidence overlap, modality agreement, temporal agreement, and transfer gaps;
  \item an HPC-oriented multi-GPU run script for the current end-to-end pipeline.
\end{itemize}

This implementation status is important for project management.
The repository is no longer only a conceptual proposal or literature review; it is now a functioning experimental scaffold that supports end-to-end dry runs and baseline experiments.

\section{Validation Status}
The current repository has been validated in three ways.
\begin{itemize}
  \item \textbf{Unit and integration testing:} the test suite covers preprocessing, retrieval, visual materialization, answerer training, oracle logic, and policy learning.
  \item \textbf{Local end-to-end dry run:} a workspace-local dry run was completed from raw TVQA-shaped annotations and synthetic videos through feature extraction, retrieval, answerer training, oracle export, policy training, and sequential evaluation.
  \item \textbf{Bug discovery before scaling:} the dry run exposed a real compatibility bug in CLIP feature extraction caused by the current \texttt{transformers} output type on the local environment; this was fixed before moving to HPC.
\end{itemize}

This is exactly the kind of validation the project should complete before large-scale experiments.
It reduces the risk of wasting GPU time on infrastructure errors and makes the transition to an HPC environment more defensible.

\section{Current Limitations}
Although the pipeline is now functional, the current implementation should still be understood as a research baseline rather than the final paper system.
The main limitations are:
\begin{itemize}
  \item the strongest implemented answerer is now a frozen multimodal scorer, but the repository still lacks a genuinely strong fine-tuned VLM-style answerer;
  \item the policy currently decides the next modality to acquire, but does not yet learn richer within-modality reranking or uncertainty-aware exploration;
  \item the local dry run used synthetic data for plumbing validation, so its metrics are not scientifically meaningful;
  \item the final claims still require experiments on real TVQA / TVQA+ data, multi-seed evaluation, robustness tests, and ablations;
  \item NExT-GQA preprocessing support is now implemented, but the benchmark still needs full end-to-end runs and evaluation support before it can substantively support the paper claims.
\end{itemize}

\section{Immediate Next Steps}
This document now captures the literature, the refined gap analysis, and the implementation status.
The next research milestones should be:

\begin{itemize}
  \item run the full pipeline on real TVQA / TVQA+ train and validation subsets;
  \item benchmark hybrid retrieval against sparse baselines under matched budgets;
  \item add a second, clearly stronger answerer tier beyond the current frozen multimodal baseline;
  \item run the new NExT-GQA preprocessing path on real data and validate the benchmark end to end;
  \item add ablations for oracle modes, stopping behavior, and policy action accuracy;
  \item add robustness experiments for noisy subtitles, missing modalities, and tighter budgets;
  \item convert the repository outputs into paper-ready tables, figures, and qualitative case studies, including transfer and evidence-overlap analysis.
\end{itemize}

\bibliographystyle{plainnat}
\bibliography{references}

\end{document}
